\clearpage\thispagestyle{empty}
\pdfbookmark[0]{Quellen zu Komet (Brief vom 1. März 1812)}{comet1812march1-sources}
\begin{center}\textbf{Quellen zu Komet (Brief vom 1. März 1812)}\end{center}
\noindent Diese KI-Übertragung ins heutige Deutsch folgt den \href{https://gdz.sub.uni-goettingen.de/id/PPN236007467}{\textit{Werken}, Band~VI} (Göttingen, 1874), S.~347--348. Gauß’ Brief ist auf den 1. März datiert und erschien als Nr.~XIX, \textit{Aus einem Schreiben des Hrn. Prof. Gauss}, im Februarheft 1812 der \textit{Monatlichen Correspondenz}, Band~XXV, S.~\href{https://api.digitale-sammlungen.de/iiif/image/v2/bsb10538614_00214/full/1600,/0/default.jpg}{206} und \href{https://api.digitale-sammlungen.de/iiif/image/v2/bsb10538614_00215/full/1600,/0/default.jpg}{207}. Das digitalisierte Exemplar stammt aus der Bayerischen Staatsbibliothek.
\par\smallskip\noindent Die Schreibweise der Beobachtungen folgt der Werkausgabe. Für 162 Mayer stehen in diesem Brief 27.0 Sekunden bei Harding und 27.5 bei Gauß. Der gesonderte Bericht vom 7. März gibt 27.0 für Gauß und 27.5 für Harding an. KI-Übertragung und KI-Satz.
\clearpage
